\section{总结}
我们描述了一种方法，它同时解决了从格点 QCD 计算部分子分布函数矩时
所面临的理论与数值挑战。
格点调节把欧几里得转动对称性破缺到超立方群，
因此格点上 twist-2 算符的重正化是确定部分子分布函数矩
$\left\langle x^{n-1} \right\rangle$ 的一大障碍。
虽然当 $n>4$ 时相对格距的幂次发散不可避免，
但对 $n=3,4$，通过选取超立方群合适的不可约表示仍可消去幂次发散；
代价是强子矩阵元的计算必须使用空间外动量，
恶化了关联函数的信噪比。
在有限流时间 $t$ 处，一旦理论裸参数与带流费米子场被重正化，
twist-2 算符的关联函数即有限。
完成连续极限后，与物理矩阵元的匹配可在短流时间处
借助理论的连续对称性完成。
使用无迹且对称化的算符保证匹配是乘法性的；
我们已在微扰论单圈水平计算出匹配系数。

研究本文的方法如何能与近期从格点 QCD 关联函数直接确定 PDF 的努力互补，
也是十分有意义的。

尽管本文提出的方法至少在原则上适用于计算部分子分布函数的任意矩，
实际上式~\eqref{eq:flowed_t2} 中带流 twist-2 算符强子矩阵元的
格点 QCD 计算受若干系统不确定度制约。

补充材料~\cite{SupplMat} 中讨论了这些系统误差，
但只有透彻的数值研究才能最终回答其适用范围，
以及这一结果是否为确定任意阶部分子分布函数矩开辟了道路。
